% Generated draft: requires mathematical and source review.
\documentclass[11pt]{article}
\usepackage[margin=25mm]{geometry}
\usepackage{fontspec}
\setmainfont{Latin Modern Roman}
\usepackage{amsmath,amssymb,mathtools,unicode-math}
\setmathfont{Latin Modern Math}
\usepackage{enumitem,needspace,xurl,hyperref}
\setlength{\parindent}{0pt}
\setlength{\parskip}{5pt}
\setlength{\emergencystretch}{3em}
\newcommand{\N}{\mathbb{N}}
\newcommand{\Z}{\mathbb{Z}}
\newcommand{\Q}{\mathbb{Q}}
\newcommand{\R}{\mathbb{R}}
\newcommand{\C}{\mathbb{C}}
\newcommand{\F}{\mathbb{F}}
\newcommand{\A}{\mathbb{A}}
\newcommand{\PP}{\mathbb{P}}
\newcommand{\E}{\mathbb{E}}
\newcommand{\eps}{\varepsilon}
\newcommand{\dd}{\,\mathrm{d}}
\begin{document}
\section*{UnsolvedMath ID: 20001907}
\textbf{Source problem code:} catalogue entry
\subsection{Definitions and mathematical statement}
The source's Conjecture \(7.2.3\), attributed to Spatzier, says that if \((M,g)\) is a closed higher-rank symmetric space of noncompact type, every metric on \(M\) with no conjugate points is a homothetic rescaling of \(g\). For a closed manifold this is read as locally symmetric, with noncompact-type symmetric universal cover. The printed statement neither passes to metrics modulo diffeomorphism nor restricts irreducibility. Both issues matter.

First, \(\varphi^*g\) has no conjugate points for every diffeomorphism \(\varphi:M\to M\), without needing to equal \(cg\) as a tensor in the original coordinates. A metric-rigidity conclusion should therefore be of the form \(h=\varphi^*(cg)\), equivalently isometry up to scale, rather than literal tensor equality.

Second, the aid gives a reducible higher-rank counterexample that survives that correction. Let \((S_1,g_1),(S_2,g_2)\) be closed hyperbolic surfaces of curvature \(-1\), and set
\[
 M=S_1\times S_2,\qquad g=g_1\oplus g_2,\qquad
 h_s=e^s g_1\oplus e^{-s}g_2\quad(s\in\R).
\]
These are nonpositively curved, hence have no conjugate points, and
\[
 \operatorname{Vol}(M,h_s)=\operatorname{Vol}(M,g).
\]
For \(s\ne0\), the two factor curvature scales differ, so \(h_s\) is not isometric to any global homothety of \(g\), even after a diffeomorphism. The aid gives the exact entropy profile
\[
 h_{\mathrm{vol}}(h_s)^2=e^{-s}+e^s=2\cosh s,
 \qquad h_{\mathrm{top}}(h_s)=h_{\mathrm{vol}}(h_s).
\]
Here volume entropy is the exponential volume-growth rate of balls in the universal cover, and topological entropy is that of the unit-speed geodesic flow. The family has fixed total volume but variable factor scales and entropy.

A viable irreducible version asks: if \((M,g)\) is a closed irreducible locally symmetric manifold of noncompact type and rank at least two, does every Riemannian \(h\) without conjugate points satisfy \(h=\varphi^*(cg)\) for some \(c>0\) and diffeomorphism \(\varphi\)? Alternatively a reducible formulation must allow factorwise scales. The supplied aid reports the irreducible no-conjugate-points version open in primary literature checked through 2026, and attributes a rigidity result under the stronger no-focal-points hypothesis to Watkins, without a full theorem statement or citation. The counterexample does not refute this corrected irreducible question.

The merged tail starts Section \(7.3\): for a closed surface \(S\) of genus at least two and a no-conjugate-points metric on \(S\times S^1\), is it a direct product? This heading and setup are preserved as a separate unfinished question. Whether direct product means the given coordinates, an isometric global product, a finite-cover product, or splitting only of the universal cover is not specified in the excerpt, and the factor-rescaling example above does not answer it.
\subsection{Short English statement}
The printed rigidity claim needs a diffeomorphism quotient and an irreducibility condition or factorwise scales. A product of hyperbolic surfaces has explicit fixed-volume, no-conjugate metrics with unequal factor scales and entropy squared \(2\cosh s\). The corrected irreducible conjecture and the appended surface-times-circle question remain separate.
\subsection{Sources}
\begin{itemize}
\item Local UnsolvedMath registry, permanent ID \(20001907\), source code \texttt{catalogue entry}; canonical collection \texttt{aim-geometry-notes.json}, notes[244]. Source item \(7.2\), visible conjecture \(7.2.3\), attributed to Spatzier.
\item AIM workshop ``Geodesics,'' Burns--Matveev report ``Open problems and questions about geodesics'': \url{https://aimath.org/pastworkshops/geodesicsproblems.pdf}. The extraction also includes the heading and introductory setup of Section \(7.3\). No PDF was opened in this drafting session.
\item The aid attributes the no-focal-points rigidity case to Watkins and qualifies the corrected irreducible no-conjugate-points literature assessment as checked through 2026. No authoritative full citation or precise theorem hypotheses are supplied. Original remarks and literature fields are empty; proposer/year registry fields are null despite the Spatzier attribution.
\item Supplied research attempt \(1\), result ``counterexample,'' run status ``refuted,'' rigor label ``proved,'' novelty confidence ``low.'' Registry status ``solved,'' classification \texttt{SOLVED-IN-LITERATURE}, updated 14 August 2026. These labels concern the printed unrestricted formulation, not the corrected conjecture.
\end{itemize}
\end{document}
